\chapter{A Concrete Naming Certificate for $\DyadicIntervalChainUp$}
\label{ch:concrete-instances-dyadic-interval-chain-namecert}
\origin{human}

Following Bishop and Bridges, the $\DyadicIntervalChainUp$ packet records a finite chain of dyadic interval cells whose inspected endpoints, nesting ledger, stream windows, regular rational readback, and terminal real-name seal are all displayed before any completion-facing consumer reads the route. It sits after the dyadic endpoint arithmetic of \autoref{ch:concrete-instances-dyadicratcore-namecert}, the finite observation windows of \autoref{ch:concrete-instances-streamname-namecert}, the regular rational readback of \autoref{ch:concrete-instances-regseqrat-namecert}, the terminal real-name boundary of \autoref{ch:concrete-instances-real-namecert}, and the finite interval-cover ledger of \autoref{ch:concrete-instances-finite-interval-cover-namecert}. The packet is a finite interval-chain route; it is not a selected limit, quotient real equality, countable choice, open-cover compactness, or ambient $\RealUp$ completeness.

\begin{definition}[Dyadic interval-chain carrier]
\label{def:dyadic-interval-chain-carrier}
A $\DyadicIntervalChainUp$ carrier is a finite $\BHist$ packet
$$
I=(D,N,W,Q,R,F,H,C,P,L)
$$
with the following displayed rows:
\begin{enumerate}
\item $D$ is the finite list of dyadic interval-cell rows, each row naming lower and upper $\DyadicRatCoreUp$ endpoints and its mesh scale.
\item $N$ is the nesting ledger, recording adjacent containment, endpoint compatibility, and the inspected width budget for the chain.
\item $W$ is the $\StreamNameUp$ window row exposing the finite stage numbers and endpoint observations used by the packet.
\item $Q$ is the $\RegSeqRatUp$ readback row attached to the displayed endpoint and width observations.
\item $R$ is the terminal $\RealUp$ seal, available only after $D,N,W,Q$ have been replayed.
\item $F$ is the finite interval-cover handoff row connecting the same cells to the finite cover surface of $\FiniteIntervalCoverUp$.
\item $H$ records componentwise $\hsame$ transport for dyadic cells, nesting rows, stream windows, regular readback, the Real seal, the finite-cover handoff, replay, provenance, and local naming.
\item $C$ records displayed $\Cont$ replay from an interval-chain request through $D,N,W,Q,R,F,H$.
\item $P$ is the $\Pkg$ provenance over $\DyadicIntervalChainUp$, $\DyadicRatCoreUp$, $\StreamNameUp$, $\RegSeqRatUp$, $\RealUp$, $\FiniteIntervalCoverUp$, $\BHist$, $\hsame$, $\Cont$, and $\NameCert$.
\item $L$ is the local $\NameCert_{\DyadicIntervalChainUp}$ row.
\end{enumerate}
The classifier compares accepted packets only through $D,N,W,Q,R,F,H,C,P,L$. It has no coordinate for a selected nested-interval limit, quotient equality of real names, countable choice, arbitrary open covers, or ambient $\RealUp$ completeness.
\end{definition}

\begin{theorem}[Dyadic interval-chain NameCert obligations]
\label{thm:dyadic-interval-chain-namecert-obligations}
For every accepted $\DyadicIntervalChainUp$ carrier
$$
I=(D,N,W,Q,R,F,H,C,P,L),
$$
the local $\NameCert_{\DyadicIntervalChainUp}$ surface consists exactly of carrier admission, dyadic endpoint exposure, adjacent nesting exactness, finite window scheduling, regular rational readback, terminal Real sealing, finite interval-cover handoff, componentwise transport, displayed replay, provenance, and local naming. Every consumer must expose the dyadic rows $D$, the nesting ledger $N$, and the finite windows $W$ before it may read $Q$, $R$, or $F$.
\end{theorem}
\begin{proof}
Project the carrier of \autoref{def:dyadic-interval-chain-carrier}. The admitted object is exactly the ten-row packet $I=(D,N,W,Q,R,F,H,C,P,L)$, so the local name can range only over the displayed dyadic cells, nesting ledger, window row, readback row, Real seal, finite-cover handoff, transport, replay, provenance, and naming rows.

The endpoint and nesting obligations are read from $D$ and $N$. The dyadic endpoint arithmetic is supplied by the $\DyadicRatCoreUp$ rows, while the nesting ledger records only adjacent containment and inspected width data.

The downstream reads are ordered by the carrier. The finite windows $W$ expose the observations used by the regular rational readback $Q$, the terminal seal $R$ is reached after those rows, and the finite-cover handoff $F$ uses the same displayed cells.

Since no other coordinate is present, the packet exports no selected limit, quotient real equality, countable choice, arbitrary open-cover compactness, or ambient completeness theorem.
\end{proof}

\begin{theorem}[Dyadic interval-chain nesting handoff]
\label{thm:dyadic-interval-chain-nesting-handoff}
Every accepted $\DyadicIntervalChainUp$ packet hands a finite nested dyadic window to later $\RealUp$ and completion-facing consumers only through the route
$$
\DyadicRatCoreUp,\ \StreamNameUp,\ \RegSeqRatUp,\ \FiniteIntervalCoverUp \longmapsto \RealUp .
$$
The handoff supplies displayed dyadic cells, adjacent containment, finite observation windows, regular rational readback, and a finite interval-cover row; it does not supply a selected limit point, quotient equality of real names, countable choice, arbitrary open-cover compactness, or ambient $\RealUp$ completeness.
\end{theorem}
\begin{proof}
Start from the accepted carrier $I=(D,N,W,Q,R,F,H,C,P,L)$. The cells $D$ and nesting ledger $N$ are the only source of interval-chain data, and both are dyadic rows.

The consumer then reads finite windows through $W$ and regular rational observations through $Q$. The terminal seal $R$ is downstream of those rows, while $F$ records the finite interval-cover handoff for the same cells.

Transport and replay through $H$ and $C$ preserve this ordering, and $P,L$ package and name the same packet. Therefore the handoff is the displayed finite route through the sibling rows and cannot be read as a selected limit, quotient equality, choice principle, open-cover compactness theorem, or ambient completeness theorem.
\end{proof}

\falsifiablePrediction{Every $\DyadicIntervalChainUp$ Real-facing read exposes dyadic interval cells, adjacent nesting, finite StreamName windows, RegSeqRat readback, a finite interval-cover handoff, componentwise transport, displayed replay, provenance, and local naming before the terminal Real seal is visible.}

\independenceWitness{dyadicratcore, streamname, regseqrat, real, finite-interval-cover: $\DyadicIntervalChainUp$ is not bijective to any listed sibling because it jointly carries finite dyadic cells, adjacent nesting, finite windows, regular rational readback, a Real seal, finite interval-cover handoff, transport, replay, provenance, and local naming.}

\closureat{DyadicIntervalChainUp}{seedStr}

\begin{closurestatus}{\DyadicIntervalChainUp}
  \constructivestory{A $\DyadicIntervalChainUp$ packet is generated from finite dyadic interval-cell rows, an adjacent nesting ledger, finite StreamName windows, RegSeqRat readback, a RealUp seal, a finite interval-cover handoff, componentwise hsame transport, displayed Cont replay, Pkg provenance, and a local NameCert row.}
  \theoryclosure{\seedClosure}
  \scopeclosed{\autoref{def:dyadic-interval-chain-carrier}, \autoref{thm:dyadic-interval-chain-namecert-obligations}, and \autoref{thm:dyadic-interval-chain-nesting-handoff} seed the finite dyadic interval-chain route from dyadic cells and nesting rows through finite windows, regular readback, finite-cover handoff, and a Real-facing seal.}
  \formalstatus{\unformalizedV}
  \bridgestatus{none}
  \notclaimed{This chapter does not close selected nested-interval limits, quotient real equality, countable choice, arbitrary open-cover compactness, ambient Real completeness, Lean verification, or bridge checking.}
  \upgradepath{Obligation closure requires checked carrier admission, dyadic endpoint exposure, adjacent nesting exactness, finite window scheduling, regular rational readback, Real seal nonescape, finite interval-cover handoff, transport, replay, provenance, and local naming for BEDC.Derived.DyadicIntervalChainUp.}
\end{closurestatus}
